● 从最基础开始                                                                                                                                                                                                                                                                                                                                                                                                        
  先理解"旋转"本身                                                                                                                                                                                            
  在三维空间里，描述一个旋转需要知道两件事：                                                                                                                                                                                                                                                                                                                                                                          
  1. 绕哪根轴转？   → 一个方向向量 n̂ = (nx, ny, nz)
  2. 转了多少度？   → 角度 θ

  比如：

  绕 Z 轴（向上）转 90°
  n̂ = (0, 0, 1)，θ = 90°

  这就是轴角表示，很直观。

  ---
  问题：为什么不直接用轴角？

  因为复合多次旋转时计算很麻烦。就像复数乘法能优雅地处理平面旋转一样，数学家发明了四元数来优雅地处理三维旋转。

  ---
  四元数是什么

  就是把轴角重新打包成4个数：

  轴角：  n̂ = (nx, ny, nz)，角度 θ

  四元数：q = [ cos(θ/2),  nx·sin(θ/2),  ny·sin(θ/2),  nz·sin(θ/2) ]
               ↑             ↑
             标量部分 w      向量部分 (x, y, z)

  例子：绕 Z 轴转 90°

  θ = 90°,  n̂ = (0, 0, 1)

  q = [ cos(45°),  0·sin(45°),  0·sin(45°),  1·sin(45°) ]
    = [ 0.707,     0,           0,            0.707       ]

  就这4个数，完整描述了这次旋转。

  ---
  为什么是 θ/2 而不是 θ？

  这是数学推导的结果。记住一个规律：

  不转（恒等旋转）：  q = [1, 0, 0, 0]   因为 cos(0/2) = 1
  转了 360°：       q = [-1, 0, 0, 0]   （等同于不转，符号相反也代表同一旋转）

  ---
  最重要的一点：为什么用四元数

  ┌─────────────────────────┬────────────────────────────────────────────────────────┐
  │          方法           │                          问题                          │
  ├─────────────────────────┼────────────────────────────────────────────────────────┤
  │ 欧拉角 (roll/pitch/yaw) │ 有万向锁：pitch=90° 时，roll 和 yaw 轴重合，丢失自由度 │
  ├─────────────────────────┼────────────────────────────────────────────────────────┤
  │ 旋转矩阵                │ 9个数，冗余，计算慢                                    │
  ├─────────────────────────┼────────────────────────────────────────────────────────┤
  │ 四元数                  │ 4个数，无奇点，复合旋转只需乘法 ✅                     │
  └─────────────────────────┴────────────────────────────────────────────────────────┘

  万向锁直观理解：

  想象用三个环嵌套来控制方向（欧拉角）
  当中间那个环转到某个角度，
  外环和内环就"对齐"了，少了一个维度
  → 某些姿态无法表示

  四元数完全避免了这个问题。

  ---
  在这个作业里的意思

  UKF 需要追踪机器人的姿态（orientation）
  姿态 = 机器人相对于世界坐标的旋转

  我们用四元数 q = [w, x, y, z] 来表示这个旋转
  初始状态：q = [1, 0, 0, 0] （没有旋转，对齐世界坐标）
  陀螺仪告诉我们每一步转了多少 → 更新 q
  最后 q → roll/pitch/yaw 用于输出

  ---
  简单说就是：四元数 = 用4个数优雅地表示三维旋转，没有数学奇点。



  